\documentclass[twocolumn]{aastex631}
\begin{document}
\title{The reionization ionizing-photon-budget shortfall for star-forming galaxies is not robust to systematics at z\~{}7 (highC systematic corner)}
\author{NebulaMind Lab (autonomous pipeline)}
\affiliation{NebulaMind Lab, \url{https://lab.nebulamind.net}}
\begin{abstract}
Reionization ionizing-photon-budget reconciliation at z\~{}7: star-forming galaxies require f\_esc=0.154 (+0.144/-0.076) to close the budget (Madau-Dickinson SFRD, log xi\_ion=25.5+/-0.15, clumping C=2-5, JWST-SFRD tail), versus indirect-proxy-inferred f\_esc=0.062 (+0.108/-0.039) from LzLCS O32/beta calibrations. Median delta(required-inferred)=+0.079 dex-frac (16-84\%: -0.036 to +0.225); 77\% of the systematic MC shows a shortfall. Conclusion: the budget CLOSES within the systematic, and the sign holds under both O32 and beta calibrations. The result is bounded by the xi\_ion x clumping x proxy-calibration systematic, not by statistics. Generated autonomously from public data (jwst) via the NebulaMind Lab runner: a bounded, reproducible, descriptive study.
\end{abstract}
\section{Introduction}
The reionization process in the early universe remains a topic of significant interest and debate among astronomers. Recent studies have highlighted potential discrepancies between the ionizing photon budget required for reionization and the observed properties of star-forming galaxies [Muñoz2024, Davies2021]. To address this issue, we revisit the ionizing-photon-budget calculation using established literature values to reconcile the cosmic star formation rate density (SFRD) with the escape fraction of ionizing photons from galaxies.
\section{Data and method}
Our approach relies on a systematics reconciliation over published literature values, avoiding direct use of survey catalog data. We adopt the Madau \& Dickinson (2014) analytic fitting function for the cosmic SFRD and utilize previously published calibrations for xi\_ion and O32/beta f\_esc proxy [LzLCS: Chisholm+22, Flury+22; Simmonds+24]. By combining these elements, we calculate the ionizing-photon-budget at z\~{}7 to assess whether star-forming galaxies can account for reionization.
\section{Result}
Our calculation reveals that star-forming galaxies require an escape fraction of f\_esc=0.154 (+0.144/-0.076) to close the ionizing-photon-budget at z\~{}7, assuming a Madau-Dickinson SFRD with log xi\_ion=25.5±0.15 and clumping factor C=2-5. In contrast, indirect-proxy-inferred escape fractions from LzLCS O32/beta calibrations yield f\_esc=0.062 (+0.108/-0.039). The median difference between the required and inferred escape fractions is +0.079 dex-frac (16-84\%: -0.036 to +0.225), with 77\% of systematic Monte Carlo simulations showing a shortfall.
\begin{figure}\centering\includegraphics[width=\columnwidth]{result.png}\caption{The reionization ionizing-photon-budget shortfall for star-forming galaxies is not robust to systematics at z\~{}7 (highC systematic corner)}\end{figure}
\section{Caveats}
It is essential to acknowledge the limitations of our approach, which relies on automated, single-selection, uncalibrated measurements. The accuracy of our result is contingent upon the validity of the adopted literature values and calibrations. Furthermore, uncertainties in xi\_ion, clumping factor, and proxy-calibration systematics introduce variability into our calculation. Therefore, while our study provides a valuable reconciliation of reionization-photon-budget discrepancies, it underscores the need for further research to refine these parameters and improve our understanding of the early universe's ionizing photon budget.
\section*{References}
\footnotesize
[Muoz2024] Muñoz et al. (2024). Reionization after JWST: a photon budget crisis?. bibcode 2024MNRAS.535L..37M \par [Park2022] Park et al. (2022). Calibrating excursion set reionization models to approximately conserve ionizing photons. bibcode 2022MNRAS.517..192P \par [Duncan2015] Duncan et al. (2015). Powering reionization: assessing the galaxy ionizing photon budget at z < 10. bibcode 2015MNRAS.451.2030D \par [Davies2021] Davies et al. (2021). The Predicament of Absorption-dominated Reionization: Increased Demands on Ionizing Sources. bibcode 2021ApJ...918L..35D \par [Madau2017] Madau (2017). Cosmic Reionization after Planck and before JWST: An Analytic Approach. bibcode 2017ApJ...851...50M

\end{document}
